\documentclass[../../main.tex]{subfiles}
\graphicspath{{\subfix{../../image/}}}

\begin{document}

\section{迹}\label{section:ch08-变分法}

下面我们讨论在 $\partial U$ 为 $C^{1}$ 的假设下，能否为函数 $u\in W^{1,p}(U)$ 沿 $\partial U$ 赋予“边界值”。若 $u\in C(\overline{U})$，则 $u$ 显然在通常意义下在 $\partial U$ 上有值。问题在于，典型的 $u\in W^{1,p}(U)$ 一般并不连续，更糟的是它只在 $U$ 中几乎处处有定义。由于 $\partial U$ 具有 $n$ 维 Lebesgue 零测度，我们无法直接给“$u$ 限制在 $\partial U$ 上”这一表达式以任何意义。迹算子的概念解决了这一问题。

\begin{theorem}[迹定理]\label{theorem:ch05-迹定理}
设$1\leqslant p<\infty$,$U\subseteq \mathbb{R}^n$是有界开集且 $\partial U$ 为 $C^{1}$。则存在有界线性算子
\begin{align*}
T:W^{1,p}(U)\to L^{p}(\partial U)
\end{align*}
使
\begin{enumerate}[(i)]
\item 若 $u\in W^{1,p}(U)\cap C(\overline{U})$，则 $Tu=u|_{\partial U}$；

\item 对每个 $u\in W^{1,p}(U)$，$\|Tu\|_{L^{p}(\partial U)}\leqslant C\|u\|_{W^{1,p}(U)}$，其中常数 $C$ 只依赖于 $p$ 与 $U$。
\end{enumerate}
我们称 $Tu$ 为 $u$ 在 $\partial U$ 上的\textbf{迹}。
\end{theorem}
\begin{proof}

1. 先设 $u\in C^{1}(\overline{U})$。如\hyperref[theorem:ch05-扩张定理]{扩张定理}证明的第一部分那样，最初也设 $x_{0}\in\partial U$ 且 $\partial U$ 在 $x_{0}$ 附近平坦，位于平面 $\{x_{n}=0\}$ 内。如前一证明中选取开球 $B$，并记$\widehat{B}$为同心、半径$\frac{r}{2}$的球。

选取 $\zeta\in C^{\infty}_{c}(B)$，在$B$中$\zeta\geqslant 0$，在$\widehat{B}$中$\zeta\equiv 1$。记 $\Gamma\triangleq \partial U\cap \widehat{B}$。令 $x'=(x_{1},\cdots,x_{n-1})\in\mathbb{R}^{n-1}=\{x_{n}=0\}$，则$\zeta$在$\partial B$上都恒为0,$\{x_n=0\}$的外法向量$n=(0,\cdots,0,-1)$,第$n$个分量为$n_{x_n}=-1$.再由\hyperref[Basis of Analytics-theorem:Green恒等式]{Green恒等式\ref{Basis of Analytics-theorem:Green恒等式-3}}可得
\begin{align*}
\int_{B^+}{\left( \zeta \left| u \right|^p \right) _{x_n}\,\mathrm{d}x}&=\int_{\partial B^+}{\zeta \left| u \right|^p\,\cdot n_{x_n}\mathrm{d}x'}=\int_{\left(\partial B\cap \left\{ x_n\geqslant 0 \right\}\right) \cup \left\{ x_n=0 \right\}}{\zeta \left| u \right|^p\,\cdot n_{x_n}\mathrm{d}x'}
\\
&=\int_{\partial B\cap \left\{ x_n\geqslant 0 \right\}}{\zeta \left| u \right|^p\,\cdot n_{x_n}\mathrm{d}x'}+\int_{\left\{ x_n=0 \right\}}{\zeta \left| u \right|^p\,\cdot n_{x_n}\mathrm{d}x'}
\\
&=\int_{\left\{ x_n=0 \right\}}{\zeta \left| u \right|^p\,\cdot n_{x_n}\mathrm{d}x'}=-\int_{\left\{ x_n=0 \right\}}{\zeta \left| u \right|^p\,\mathrm{d}x'}.
\end{align*}
从而
\begin{align}
\int_{\Gamma}{|u|^p\,\mathrm{d}x'}&=\int_{\Gamma}{\zeta |u|^p\,\mathrm{d}x'}\leqslant \int_{\{x_n=0\}}{\zeta |u|^p\,\mathrm{d}x'}=-\int_{B^+}{(\zeta |u|^p)_{x_n}\,\mathrm{d}x}\nonumber
\\
&=-\int_{B^+}{|u|^p\zeta _{x_n}}+p|u|^{p-1}\left( \sgn u \right) u_{x_n}\zeta \,\mathrm{d}x\nonumber
\\
&\leqslant -\int_{B^+}{|u|^p\zeta _{x_n}}+p|u|^{p-1}\left| u_{x_n}\, \right|\mathrm{d}x
\nonumber\\
&\overset{\text{\hyperref[Basis of Analytics-theorem:Young不等式]{Young不等式}}}{\leqslant}C\int_{B^+}{|u|^p\,\mathrm{d}x}+\int_{B^+}{\left( \left( p-1 \right) \left| u \right|^p+\left| u_{x_n} \right|^p \right) \,\mathrm{d}x}\nonumber
\\
&\leqslant C\int_{B^+}{|u|^p}+|Du|^p\,\mathrm{d}x,
\label{align:ch06-内部H2估计}
\end{align}
其中$C$为某一固定常数.进而
\begin{align*}
\left\| u \right\| _{L^p\left( \Gamma \right)}\leqslant C\left\| u \right\| _{W^{1,p}\left( B^+ \right)}.
\end{align*}

2. 若 $x_{0}\in\partial U$，但 $\partial U$ 在 $x_{0}$ 附近不平坦，我们利用微分同胚拉直边界。由 $\partial U$ 是 $C^{1}$，存在开集 $W\subseteq\mathbb{R}^{n}$ 包含 $x_{0}$，以及 $C^{1}$ 微分同胚
\begin{align*}
\boldsymbol{\Phi}: W \to B(0,r),\qquad \boldsymbol{\Psi}=\boldsymbol{\Phi}^{-1},
\end{align*}
使得
\begin{align*}
\boldsymbol{\Phi}(W\cap U)=B^{+}(0,r),\qquad
\boldsymbol{\Phi}(W\cap\partial U)=B(0,r)\cap\{x_{n}=0\}.
\end{align*}
记 $\Gamma=\partial U\cap W$。先设 $u\in C^{1}(\overline{U})$，定义
\begin{align*}
v(y)\triangleq u(\boldsymbol{\Psi}(y)),\qquad y\in B^{+}(0,r).
\end{align*}
则 $v\in C^{1}(\overline{B^{+}(0,r)})$。由步骤 1 中平坦情形的估计 \eqref{align:ch06-内部H2估计}，对 $\Gamma'=\boldsymbol{\Phi}(\Gamma)\subseteq B(0,r)\cap\{x_{n}=0\}$ 有
\begin{align}\label{eq:ch05-迹定理-拉直估计}
\int_{\Gamma'}|v(y)|^{p}\,\mathrm{d}S_{y}
\leqslant C\int_{B^{+}(0,r)}(|v|^{p}+|Dv|^{p})\,\mathrm{d}y,
\end{align}
其中 $C$ 只依赖于 $p$ 与 $n$。

现作变量代换 $y=\boldsymbol{\Phi}(x)$。由于 $\boldsymbol{\Phi}$ 是 $C^{1}$ 微分同胚，其限制在边界 $\partial U$ 上也是 $C^{1}$ 的，因此存在只依赖于 $\boldsymbol{\Phi}$ 的常数 $C_{1}>0$，使得面积元满足
\begin{align*}
\mathrm{d}S_{y}\leqslant C_{1}\,\mathrm{d}S_{x},
\end{align*}
从而
\begin{align}\label{eq:ch05-迹定理-面积元}
\int_{\Gamma}|u(x)|^{p}\,\mathrm{d}S_{x}
\leqslant C_{1}\int_{\Gamma'}|v(y)|^{p}\,\mathrm{d}S_{y}.
\end{align}
另一方面，由 $y=\boldsymbol{\Phi}(x)$ 得 $\mathrm{d}y=|\det D\boldsymbol{\Phi}(x)|\,\mathrm{d}x$，且由$\boldsymbol{\Phi}$是$C^1$的和链式法则知,存在常数 $C_{2}>0$ 使得
\begin{align*}
|\det D\boldsymbol{\Phi}(x)|\leqslant C_{2},\qquad
|Dv(y)|=|Du(x)D\boldsymbol{\Psi}(y)|\leqslant C_{2}|Du(x)|.
\end{align*}
于是
\begin{align}\label{eq:ch05-迹定理-体积元}
\int_{B^{+}(0,r)}(|v|^{p}+|Dv|^{p})\,\mathrm{d}y
\leqslant C_{3}\int_{W\cap U}(|u|^{p}+|Du|^{p})\,\mathrm{d}x
\leqslant C_{3}\int_{U}(|u|^{p}+|Du|^{p})\,\mathrm{d}x,
\end{align}
其中 $C_{3}$ 只依赖于 $n,p$ 与 $\boldsymbol{\Phi}$。将 \eqref{eq:ch05-迹定理-面积元} 和 \eqref{eq:ch05-迹定理-体积元} 代入 \eqref{eq:ch05-迹定理-拉直估计}，得到
\begin{align*}
\int_{\Gamma}|u|^{p}\,\mathrm{d}S
\leqslant C\int_{U}(|u|^{p}+|Du|^{p})\,\mathrm{d}x.
\end{align*}
因此
\begin{align*}
\|u\|_{L^{p}(\Gamma)}\leqslant C\|u\|_{W^{1,p}(U)},
\end{align*}
其中常数 $C$ 只依赖于 $p$，$U$ 与拉直映射 $\boldsymbol{\Phi}$，但不依赖于 $u$。而拉直映射 $\boldsymbol{\Phi}$只依赖于$U$，故常数 $C$ 实际上只依赖于 $p$，$U$.

3. 由于 $\partial U$ 紧，存在有限个点 $x_{i}^{0}\in\partial U$和对应邻域$B_i$及开子集 $\Gamma_{i}=\partial U\cap B_i\subset\partial U$（$i=1,\cdots,N$），使 $\partial U=\bigcup_{i=1}^{N}\Gamma_{i}$ 且
\begin{align*}
\|u\|_{L^{p}(\Gamma_{i})}\leqslant C\|u\|_{W^{1,p}(U)}\quad (i=1,\cdots,N).
\end{align*}
因而若写
\begin{align*}
Tu\triangleq u|_{\partial U},
\end{align*}
则
\begin{align}\label{eq::FJjioFJIOWwwqqqq}
\|Tu\|_{L^{p}(\partial U)}\leqslant C\|u\|_{W^{1,p}(U)}
\end{align}
其中常数 $C$ 与 $u$ 无关。

4. 不等式\eqref{eq::FJjioFJIOWwwqqqq}对 $u\in C^{1}(\overline{U})$ 成立。现设 $u\in W^{1,p}(U)$，则由\cref{theorem:ch05-边界光滑逼近}知存在函数 $u_{m}\in C^{\infty}(\overline{U})$ 在 $W^{1,p}(U)$ 中收敛于 $u$。由上述不等式\eqref{eq::FJjioFJIOWwwqqqq}有
\begin{align}
\|Tu_{m}-Tu_{l}\|_{L^{p}(\partial U)}\leqslant C\|u_{m}-u_{l}\|_{W^{1,p}(U)};
\end{align}
故 $\{Tu_{m}\}_{m=1}^{\infty}$ 是 $L^{p}(\partial U)$ 中的 Cauchy 序列。而$L^{p}(\partial U)$是完备的,我们定义
\begin{align*}
Tu\triangleq\lim_{m\to\infty}Tu_{m},
\end{align*}
此定义不依赖于逼近 $u$ 的光滑函数的具体选取。事实上，设 \(\{u_m\}_{m=1}^\infty\) 与 \(\{v_m\}_{m=1}^\infty\) 均为 \(C^\infty(\overline U)\) 中的序列，并且在 \(W^{1,p}(U)\) 中收敛于 \(u\)。由第 3 步的不等式 \eqref{eq::FJjioFJIOWwwqqqq}，对任意 \(m\) 有
\begin{align*}
\|Tu_m-Tv_m\|_{L^p(\partial U)}
\leqslant C\|u_m-v_m\|_{W^{1,p}(U)}.
\end{align*}
由于
\begin{align*}
\|u_m-v_m\|_{W^{1,p}(U)}
\leqslant \|u_m-u\|_{W^{1,p}(U)}+\|u-v_m\|_{W^{1,p}(U)},
\end{align*}
且右端当 \(m\to\infty\) 时趋于 \(0\)，故
\begin{align*}
\|Tu_m-Tv_m\|_{L^p(\partial U)}\to 0\quad (m\to\infty).
\end{align*}
另一方面，由 \eqref{eq::FJjioFJIOWwwqqqq}式可得 \(\{Tu_m\}\) 与 \(\{Tv_m\}\) 都是 \(L^p(\partial U)\) 中的 Cauchy 列，因而分别收敛于某 \(g,h\in L^p(\partial U)\)。于是
\begin{align*}
\|g-h\|_{L^p(\partial U)}
=\lim\limits_{m\to\infty}\|Tu_m-Tv_m\|_{L^p(\partial U)}
=0.
\end{align*}
所以 \(g=h\) 在 \(L^p(\partial U)\) 中成立。因此极限 \(Tu\triangleq\lim\limits_{m\to\infty}Tu_m\) 不依赖于逼近序列 \(\{u_m\}_{m=1}^\infty\) 的具体选取。

最后若 $u\in W^{1,p}(U)\cap C(\overline{U})$，注意此时\cref{theorem:ch05-边界光滑逼近}证明中构造的函数$u_{m}\in C^{\infty}(\overline{U})$在$\overline{U}$上一致收敛于$u$,因而
\begin{align*}
Tu=\underset{m\rightarrow \infty}{\lim}Tu_m=\underset{m\rightarrow \infty}{\lim}u_m|_{\partial U}=u|_{\partial U}.
\end{align*}
\end{proof}

\begin{theorem}[$W^{1,p}$ 中的零迹函数]\label{theorem:ch05-零迹函数}
设$U\subseteq \mathbb{R}^n$是有界开集且 $\partial U$ 为 $C^{1}$,$1\leqslant p<\infty$,$u\in W^{1,p}(U)$。则
\begin{align*}
\text{在}\partial U\text{上}\text{有}Tu=0\text{当且仅当}u\in W^{1,p}_{0}(U).
\end{align*}
\end{theorem}
\begin{remark}
感觉$p=\infty$的时候似乎也成立，充分性同样证明,必要性应该能利用\cref{theorem:ch05-W1infty刻画}证明.

因为针对 \(p<\infty\) 的标准证明依赖于 \(C^\infty(\overline{U})\) 在 \(W^{1,p}(U)\) 中的稠密性，而这一稠密性在 \(W^{1,\infty}(U)\) 中并不成立（例如，\(u(x)=|x|\) 在 \((-1,1)\) 上属于 \(W^{1,\infty}\)，但无法被 \(C^1\) 函数在 \(W^{1,\infty}\) 范数下逼近）。因此，必要性的证明不能直接把 \(p<\infty\) 的证明照搬过来。
\end{remark}
\begin{proof}

{\heiti 充分性:} 设 $u\in W^{1,p}_{0}(U)$。由定义，存在函数列 $u_{m}\in C^{\infty}_{c}(U)$ 使得
\begin{align*}
u_{m}\to u\quad\text{于}\ W^{1,p}(U).
\end{align*}
因为 $u_{m}$ 在 $U$ 内有紧支集，故 $Tu_{m}=0$ 于 $\partial U$。又由\hyperref[theorem:ch05-迹定理]{迹定理}，迹算子
\begin{align*}
T:W^{1,p}(U)\to L^{p}(\partial U)
\end{align*}
是有界线性算子，因此
\begin{align*}
Tu=\lim_{m\to\infty}Tu_{m}=0\quad\text{于}\ L^{p}(\partial U).
\end{align*}
故 $Tu=0$ 于 $\partial U$。

{\heiti 必要性:} 设 $u\in W^{1,p}(U)$ 且
\begin{align}\label{eq:ch05-零迹-假设}
Tu=0\quad\text{于}\ \partial U.
\end{align}
下面证明 $u\in W^{1,p}_{0}(U)$。

1.
设 $v\in W^{1,p}(\mathbb{R}^{n}_{+})$，在 $\overline{\mathbb{R}^{n}_{+}}$ 中有紧支集，且
\begin{align}\label{eq:ch05-零迹-半空间假设}
Tv=0\quad\text{于}\ \partial\mathbb{R}^{n}_{+}=\mathbb{R}^{n-1}.
\end{align}
证明 $v\in W^{1,p}_{0}(\mathbb{R}^{n}_{+})$。

由\hyperref[theorem:ch05-边界光滑逼近]{边界光滑函数的逼近定理}以及迹算子 $T$ 的连续性，存在函数列 $v_{m}\in C^{\infty}(\overline{\mathbb{R}^{n}_{+}})$ 使得
\begin{align}\label{eq:ch05-零迹-半空间逼近}
v_{m}\to v\quad\text{于}\ W^{1,p}(\mathbb{R}^{n}_{+}),
\end{align}
并且
\begin{align}\label{eq:ch05-零迹-半空间迹逼近}
Tv_{m}=v_{m}|_{\mathbb{R}^{n-1}}\to Tv=0\quad\text{于}\ L^{p}(\mathbb{R}^{n-1}).
\end{align}

对每个 $m$，由微积分基本定理，对任意 $x=(x',x_{n})\in\mathbb{R}^{n}_{+}$，有
\begin{align*}
|v_{m}(x',x_{n})|\leqslant |v_{m}(x',0)|+\int_{0}^{x_{n}}|\partial_{n}v_{m}(x',t)|\,\mathrm{d}t.
\end{align*}
两边取 $p$ 次方，利用不等式 $(a+b)^{p}\leqslant 2^{p-1}(a^{p}+b^{p})$，再对 $x'$ 在 $\mathbb{R}^{n-1}$ 上积分，并用\hyperref[Real Analysis-theorem:实变函数--Hölder不等式]{Hölder不等式}，得到
\begin{align*}
\int_{\mathbb{R}^{n-1}}|v_{m}(x',x_{n})|^{p}\,\mathrm{d}x'
&\leqslant 2^{p-1}\int_{\mathbb{R}^{n-1}}|v_{m}(x',0)|^{p}\,\mathrm{d}x' +2^{p-1}\int_{\mathbb{R}^{n-1}}\left(\int_{0}^{x_{n}}|\partial_{n}v_{m}(x',t)|\,\mathrm{d}t\right)^{p}\mathrm{d}x'\\
&\leqslant C\left(\int_{\mathbb{R}^{n-1}}|v_{m}(x',0)|^{p}\,\mathrm{d}x'
+x_{n}^{p-1}\int_{0}^{x_{n}}\int_{\mathbb{R}^{n-1}}|Dv_{m}(x',t)|^{p}\,\mathrm{d}x'\,\mathrm{d}t\right),
\end{align*}
其中常数 $C$ 只依赖于 $p$。令 $m\to\infty$，由 \eqref{eq:ch05-零迹-半空间逼近} 和 \eqref{eq:ch05-零迹-半空间迹逼近}，推得对几乎处处的 $x_{n}>0$，有
\begin{align}\label{eq:ch05-零迹-半空间估计}
\int_{\mathbb{R}^{n-1}}|v(x',x_{n})|^{p}\,\mathrm{d}x'
\leqslant Cx_{n}^{p-1}\int_{0}^{x_{n}}\int_{\mathbb{R}^{n-1}}|Dv(x',t)|^{p}\,\mathrm{d}x'\,\mathrm{d}t.
\end{align}

取 $\zeta\in C^{\infty}(\mathbb{R}_{+})$ 满足
\begin{align*}
0\leqslant \zeta\leqslant 1,\qquad
\zeta\equiv 1\ \text{于}\ [0,1],\qquad
\zeta\equiv 0\ \text{于}\ [2,\infty).
\end{align*}
令
\begin{align*}
\zeta_{m}(x)\triangleq \zeta(mx_{n})\quad (x\in\mathbb{R}^{n}_{+}),\qquad
w_{m}\triangleq v(1-\zeta_{m}).
\end{align*}
则
\begin{align*}
\partial_{n}w_{m}=\partial_{n}v(1-\zeta_{m})-mv\zeta'(mx_{n}),\qquad
D_{x'}w_{m}=D_{x'}v(1-\zeta_{m}).
\end{align*}
于是
\begin{align*}
Dw_{m}-Dv
&=\big(D_{x'}w_{m}-\ D_{x'}v,\ \partial_{n}w_{m}-\partial_{n}v\big)\\
&=\big(-D_{x'}v\,\zeta_{m},\ -\partial_{n}v\,\zeta_{m}-m v\zeta'(mx_{n})\big)\\
&=-\zeta_{m}Dv - m v\zeta'(mx_{n})e_{n},
\end{align*}
其中 $e_{n}=(0,\cdots,0,1)$。因此由三角不等式和 $(a+b)^{p}\leqslant 2^{p-1}(a^{p}+b^{p})$ 得
\begin{align*}
|Dw_{m}-Dv|^{p}
\leqslant \big(|\zeta_{m}||Dv| + m|v||\zeta'(mx_{n})|\big)^{p}\leqslant 2^{p-1}\Big(|\zeta_{m}|^{p}|Dv|^{p} + m^{p}|v|^{p}|\zeta'(mx_{n})|^{p}\Big).
\end{align*}
两边在 $\mathbb{R}^{n}_{+}$ 上积分，得到
\begin{align*}
\int_{\mathbb{R}^{n}_{+}}|Dw_{m}-Dv|^{p}\,\mathrm{d}x
\leqslant 2^{p-1}\int_{\mathbb{R}^{n}_{+}}|\zeta_{m}|^{p}|Dv|^{p}\,\mathrm{d}x + 2^{p-1}m^{p}\int_{\mathbb{R}^{n}_{+}}|v|^{p}|\zeta'(mx_{n})|^{p}\,\mathrm{d}x.
\end{align*}
由 $\zeta$ 的选取，$\zeta'\equiv 0$ 于 $[0,1]\cup[2,\infty)$，故 $\zeta'(mx_{n})$ 仅在 $1< mx_{n}<2$，即$\frac{1}{m}< x_{n}<\frac{2}{m}$时可能非零。又因为 $|\zeta'|\leqslant M$ 对某个只依赖于 $\zeta$ 的常数 $M$ 成立，所以
\begin{align*}
\int_{\mathbb{R}^{n}_{+}}|v|^{p}|\zeta'(mx_{n})|^{p}\,\mathrm{d}x
&\leqslant M^{p}\int_{0}^{\frac{2}{m}}\int_{\mathbb{R}^{n-1}}|v(x',x_{n})|^{p}\,\mathrm{d}x'\,\mathrm{d}x_{n}.
\end{align*}
将上述两式合并，并取 $C=\max\{2^{p-1},\,2^{p-1}M^{p}\}$，即得
\begin{align}\label{eq:ch05-零迹-估计AB}
\int_{\mathbb{R}^{n}_{+}}|Dw_{m}-Dv|^{p}\,\mathrm{d}x
\leqslant C\int_{\mathbb{R}^{n}_{+}}|\zeta_{m}|^{p}|Dv|^{p}\,\mathrm{d}x
+Cm^{p}\int_{0}^{\frac{2}{m}}\int_{\mathbb{R}^{n-1}}|v|^{p}\,\mathrm{d}x'\,\mathrm{d}t
\triangleq A+B.
\end{align}
当 $m\to\infty$ 时，$A\to 0$，因为 $\zeta_{m}$ 仅在 $0\leqslant x_{n}\leqslant \frac{2}{m}$ 上非零，而该区域测度趋于零。为估计 $B$，利用 \eqref{eq:ch05-零迹-半空间估计} 得
\begin{align*}
B
&\leqslant Cm^{p}\int_{0}^{\frac{2}{m}}\left(\int_{\mathbb{R}^{n-1}}|v(x',t)|^{p}\,\mathrm{d}x'\right)\mathrm{d}t\\
&\leqslant Cm^{p}\int_{0}^{\frac{2}{m}}\left(Ct^{p-1}\int_{0}^{t}\int_{\mathbb{R}^{n-1}}|Dv(x',s)|^{p}\,\mathrm{d}x'\,\mathrm{d}s\right)\mathrm{d}t\\
&\leqslant Cm^{p}\left(\int_{0}^{\frac{2}{m}}t^{p-1}\,\mathrm{d}t\right)
\left(\int_{0}^{\frac{2}{m}}\int_{\mathbb{R}^{n-1}}|Dv(x',s)|^{p}\,\mathrm{d}x'\,\mathrm{d}s\right)\\
&\leqslant C'\int_{0}^{\frac{2}{m}}\int_{\mathbb{R}^{n-1}}|Dv(x',s)|^{p}\,\mathrm{d}x'\,\mathrm{d}s\to 0,\quad m\to\infty.
\end{align*}
于是由 \eqref{eq:ch05-零迹-估计AB} 得 $Dw_{m}\to Dv$ 于 $L^{p}(\mathbb{R}^{n}_{+})$。又显然 $w_{m}\to v$ 于 $L^{p}(\mathbb{R}^{n}_{+})$，故
\begin{align*}
w_{m}\to v\quad\text{于}\ W^{1,p}(\mathbb{R}^{n}_{+}).
\end{align*}
但当 $0<x_{n}<\frac{1}{m}$ 时，$\zeta_{m}=1$，从而 $w_{m}=0$。因此 $w_{m}$ 的支集与边界 $\{x_{n}=0\}$ 的距离至少为 $\frac{1}{m}$。又$v$在 $\overline{\mathbb{R}^{n}_{+}}$ 中有紧支集,故由\hyperref[theorem:ch05-局部光滑逼近]{光滑函数的局部逼近定理},对 $w_{m}$ 进行磨光，可得到函数 $v_{m}\in C^{\infty}_{c}(\mathbb{R}^{n}_{+})$，使得
\begin{align*}
v_{m}\to v\quad\text{于}\ W^{1,p}(\mathbb{R}^{n}_{+}).
\end{align*}
故 $v\in W^{1,p}_{0}(\mathbb{R}^{n}_{+})$。

2.由于 $\partial U$ 是紧的 $C^{1}$ 边界，存在有限个开集 $W_{i}$（$i=1,\cdots,N$）覆盖 $\partial U$，以及相应的 $C^{1}$ 微分同胚
\begin{align*}
\boldsymbol{\Phi}_{i}:W_{i}\to B(0,r_{i}),\qquad \boldsymbol{\Psi}_{i}=\boldsymbol{\Phi}_{i}^{-1},
\end{align*}
使得
\begin{align}\label{eq:ch05-零迹-拉直}
\boldsymbol{\Phi}_{i}(W_{i}\cap U)=B^{+}(0,r_{i}),\qquad
\boldsymbol{\Phi}_{i}(W_{i}\cap\partial U)=B(0,r_{i})\cap\{x_{n}=0\},
\end{align}
其中 $B^{+}(0,r_{i})=B(0,r_{i})\cap\{x_{n}>0\}$。再取 $W_{0}\subset\subset U$ 使得
\begin{align*}
U\subseteq \bigcup_{i=0}^{N}W_{i}.
\end{align*}
由\hyperref[lemma:ch05-光滑单位分解-p7R3]{光滑单位分解引理}，存在从属于覆盖 $\{W_{i}\}_{i=0}^{N}$ 的光滑单位分解 $\{\zeta_{i}\}_{i=0}^{N}$，即
\begin{align*}
\zeta_{i}\in C_{c}^{\infty}(W_{i}),\qquad 0\leqslant \zeta_{i}\leqslant 1,\qquad
\sum_{i=0}^{N}\zeta_{i}=1\quad\text{于}\ U.
\end{align*}
于是
\begin{align}\label{eq:ch05-零迹-分解}
u=\sum_{i=0}^{N}\zeta_{i}u.
\end{align}
对于 $i=0$，$\zeta_{0}u$ 具有紧支集在 $U$ 内，由\hyperref[theorem:ch05-局部光滑逼近]{光滑函数的局部逼近定理}知,存在函数 $f_{m}\in C^{\infty}_{c}(U)$，使得
\begin{align*}
f_{m}\to \zeta_{0}u\quad\text{于}\ W^{1,p}(U).
\end{align*}
故 $\zeta_{0}u\in W^{1,p}_{0}(U)$。

对于 $i=1,\cdots,N$，$\zeta_{i}u$ 的支集在 $W_{i}$ 内。令
\begin{align*}
v_{i}(y)\triangleq (\zeta_{i}u)(\boldsymbol{\Psi}_{i}(y)),\qquad y\in B^{+}(0,r_{i}).
\end{align*}
则 $v_{i}\in W^{1,p}(B^{+}(0,r_{i}))$，且 $v_{i}$ 在 $B^{+}(0,r_{i})$ 中有紧支集。由\eqref{eq:ch05-零迹-假设} 和 \eqref{eq:ch05-零迹-拉直}，有
\begin{align*}
Tv_{i}=0\quad\text{于}\ B(0,r_{i})\cap\{x_{n}=0\}.
\end{align*}
将 $v_{i}$ 零延拓到整个 $\mathbb{R}^{n}_{+}$，仍记为 $v_{i}$。则 $v_{i}\in W^{1,p}(\mathbb{R}^{n}_{+})$，在 $\overline{\mathbb{R}^{n}_{+}}$ 中有紧支集，且 $Tv_{i}=0$ 于 $\partial\mathbb{R}^{n}_{+}$。由第一步的半空间情形，得
\begin{align*}
v_{i}\in W^{1,p}_{0}(\mathbb{R}^{n}_{+}),
\end{align*}
从而 $v_{i}\in W^{1,p}_{0}(B^{+}(0,r_{i}))$。通过 $\boldsymbol{\Psi}_{i}$ 拉回，得到
\begin{align*}
\zeta_{i}u\in W^{1,p}_{0}(W_{i}\cap U).
\end{align*}
将 $\zeta_{i}u$ 零延拓到整个 $U$ 上，仍记为 $\zeta_{i}u$，则 $\zeta_{i}u\in W^{1,p}_{0}(U)$。

最后，由 \eqref{eq:ch05-零迹-分解} 及 $W^{1,p}_{0}(U)$ 是线性空间，有限和仍属于 $W^{1,p}_{0}(U)$，故
\begin{align*}
u=\sum_{i=0}^{N}\zeta_{i}u\in W^{1,p}_{0}(U).
\end{align*}
\end{proof}

\begin{examples}
设 \(U = (0,1)\)，其边界为 \(\partial U = \{0, 1\}\)。此时迹算子 \(T\) 是从 \(W^{1,2}(U)\) 到 \(L^2(\partial U)\) 的连续线性映射。若取 \(u(x) = x^2\)，它是 \(W^{1,2}(U)\) 中的函数，并且因为它是连续的，它的迹就是直接取边界值，即
\begin{align*}
Tu = (u(0), u(1)) = (0, 1).
\end{align*}
这说明 \(u\) 在边界上不为零，因此 \(u \notin W_0^{1,2}(0,1)\)。若取 \(u(x) = x(1-x)\)，它的迹为
\begin{align*}
Tu = (u(0), u(1)) = (0, 0).
\end{align*}
因此 \(u\) 属于 \(W_0^{1,2}(0,1)\)。这个例子表明，\(W_0^{1,2}(U)\) 的核空间确实精确对应着在边界上取值为零的函数，而迹算子 \(T\) 则通过连续延拓将“边界取值”这一概念从经典函数推广到了整个Sobolev空间。
\end{examples}

\begin{examples}
对于经典的Dirichlet问题
\begin{align*}
-\Delta u = f \text{ 于 } U \text{ 中}, \qquad u = g \text{ 于 } \partial U \text{ 上}.
\end{align*}
如果直接在 \(W^{1,2}(U)\) 中寻找弱解，由于函数仅可积，逐点满足 \(u = g\) 在边界上毫无意义。借助迹算子，我们可以将边界条件弱化为 \(Tu = g\)，等价地表述为 \(T(u-g) = 0\)，即 \(u - g \in W_0^{1,2}(U)\)。
将方程两端乘以测试函数 \(v \in W_0^{1,2}(U)\) 并分部积分，得到弱形式(见\cref{Functional Analysis-definition:hjojwfeiowjojfiojwiowej})
\begin{align*}
\int_U \nabla u \cdot \nabla v \, \mathrm{d}x = \int_U f v \, \mathrm{d}x.
\end{align*}
这里的测试函数 \(v\) 取在 \(W_0^{1,2}(U)\) 中是至关重要的，因为分部积分产生的边界项为 \(\int_{\partial U} \frac{\partial u}{\partial \nu} v \, \mathrm{d}S\)。由于 \(v \in W_0^{1,2}(U)\)，其迹 \(Tv = 0\)，因此这个边界项自动消失。这样一来，弱解的定义既能严格表达“在边界上取给定值”这一物理条件，又允许我们使用\hyperref[Functional Analysis-theorem:泛函分析--Lax-Milgram定理]{Lax-Milgram定理}等泛函分析工具放在更弱的$W_0^{1,2}(U)$中来证明解的存在唯一性。
\end{examples}



\end{document}
